\section{The original canonical cusp connection over arbitrary coefficient algebras}
\label{sec:cnc-v47}

We construct the connection required by the original predicate
\(\mathrm{IsNab}\).  In particular, the connection hypothesis retained in
ICL11 is supplied on each original moduli object.  The construction uses
the given universal abelian scheme and its line bundle; it does not
choose a different moduli space, complex embedding, Hermitian form,
metric, measure, or tuple phase.  Principal analytic charts and the
finite theta system remain separate producers.

Fix a CM number field \(E\), its original host \(H\), a datum
\(D_w\), a level \(l\), a moduli object \(X_o:\mathrm{Moduli}(D_w,l)\),
and the original coefficient field \(K=c.F\).  Write
\[
 \begin{gathered}
 O=\mathcal O_E,\qquad A=X_o.\mathrm{univ}.t.G,
 \qquad L=X_o.\mathrm{univ}.t.PP.L,\\
 I=\coprod_{\tau:E\hookrightarrow\mathbb C}
            \{0,\ldots,D_w.\mathrm{mult}(\tau)-1\}.
 \end{gathered}
                                                               \tag{CNC.1}
\]
The involution on \(O\) is \(a^*=H.\mathrm{bar}(a)\).  It is an
involution because its image in \(E\) is complex conjugation.  The
element \(w_0\) is the original \(H.w_0\).
For a commutative \(K\)-algebra \(S\), no finiteness, reducedness,
locality, or domain hypothesis is imposed.  Put
\[
 J_m(S,n)=S[t_{b,i}:1\leq b\leq n,\ 1\leq i\leq d]
                    /(t)^{m+1}.
\]
All infinite series used below are formal series.  Their substitutions,
exponentials, logarithms, and integrals are computed coefficientwise;
at each fixed degree every calculation takes place in a finite jet
algebra.  There is no analytic convergence assertion about \(S\).

Two general geometric foundations enter.  A smooth morphism with a
section has locally power-series formal coordinates along that section;
the conormal graded algebra is its symmetric algebra.  Also, a proper
flat morphism of finite presentation with geometrically reduced connected
fibres has universally only base functions.  Precise primary references
are the Stacks Project, tags
\href{https://stacks.math.columbia.edu/tag/05D9}{05D9} and
\href{https://stacks.math.columbia.edu/tag/0E0L}{0E0L}, respectively.
The passage from local coordinates to the globally framed coefficient
algebra, the normalization of all line generators, and the formal gauge
equations are proved here.  The supplied cubical, adjoint, and primitive
line-bundle identifications are fields of the original \(\mathrm{SAB}\)
object, not additional existence hypotheses.

% BEGIN CNC1
\paragraph{CNC1. The given family and uniqueness of normalized line isomorphisms.}
\label{claim:CNC1}
The universal group \(A\to X_o.X\) and each of its finite powers are
abelian schemes.  After every base change \(T\to X_o.X\),
\[
 \mathcal O_T\ \xrightarrow{\ \sim\ }\ (A_T^r\to T)_*
                       \mathcal O_{A_T^r}\qquad(r\geq0).       \tag{CNC.2}
\]
Consequently two isomorphisms between the same two line bundles on
\(A_T^r\) agree if their ratio is one on the zero section.  This assertion
holds for nonreduced \(T\).
\begin{proof}
The original \(\mathrm{SAB}.\mathrm{semi}\) supplies smoothness,
commutativity and geometrically connected fibres.  The field
\(\mathrm{Obj}.\mathrm{proper}\), applied to \(X_o.\mathrm{univ}\),
supplies properness.  These are exactly the four fields of the original
\(\mathrm{IsAbelianScheme}\), as is also visible in the source proof of
\texttt{ObjDim.isAbelianScheme}.  Smoothness gives flatness, local finite
presentation and geometrically reduced fibres.  Properness supplies
quasi-compactness and quasi-separatedness, so the morphism is of finite
presentation.  These properties persist under base change and finite
products; the product of the geometrically connected smooth fibres is
geometrically connected and reduced.  Apply the universal base-functions
theorem just stated.  For \(r=0\) the map is the identity of \(T\).

The ratio of two line-bundle isomorphisms is a global invertible
function: this follows in a local basis and the ratios agree on overlaps.
By (CNC.2) it is pulled back from a unique invertible function on \(T\).
The zero section is a left inverse of this pullback, so ratio one there
forces ratio one everywhere.  This also proves uniqueness of a line
isomorphism after its zero-section normalization.  The argument uses
actual functions and their pullbacks, not tests on geometric points.
\end{proof}
% END CNC1

% BEGIN CNC2
\paragraph{CNC2. The original eigenframes form the whole tangent basis.}
\label{claim:CNC2}
For every original framed point \(y\in\mathrm{FrPt}(X_o,c,S)\), its
vectors \(y.u_{\tau,i}\), indexed by \(I\), form an \(S\)-basis of
the whole tangent module of \(A_S\) at zero.
\begin{proof}
The definition of \(\mathrm{Coeff}\) puts every embedding of \(E\)
inside \(K\).  Hence
\[
 E\otimes_{\mathbb Q}K\ \longrightarrow\prod_{\tau:E\hookrightarrow K}K,
       \qquad a\otimes b\longmapsto(\tau(a)b)_\tau             \tag{CNC.3}
\]
is an isomorphism.  Here is an explicit verification sufficient for
arbitrary \(S\).  Choose an integral basis \(b_1,\ldots,b_e\) of
\(O\).  The matrix \((\tau(b_j))_{\tau,j}\) is invertible over
\(K\): its squared determinant is the nonzero discriminant, after the
chosen embeddings are ordered.  Its inverse expresses every coordinate
idempotent of the product in the basis \(b_j\).  It also proves
multiplicativity and bijectivity of (CNC.3), by checking on the displayed
coordinates.

The differential of the given \(O\)-action extends uniquely
\(\mathbb Q\)-linearly to \(E\).  Combining it with the scalar
\(K\)-action gives an action of the left side of (CNC.3) on the tangent
module.  Its coordinate idempotents \(p_\tau\) are mutually orthogonal
and sum to one.  Their images are exactly
\[
 M_\tau=\{v:\ a v=\tau(a)v\text{ for every }a\in O\}.
\]
Indeed the coordinate relation proves the inclusion into \(M_\tau\);
conversely the inverse embedding matrix shows that \(p_\sigma v\)
is zero for \(\sigma\ne\tau\) and equals \(v\) for
\(\sigma=\tau\).  Thus the whole tangent module is
\(\bigoplus_\tau M_\tau\), over the given \(S\), without making
any assertion about integral eigenspaces at ramified primes.
The fields \(y.\mathrm{tan}\), \(y.\mathrm{eig}\), and
\(y.\mathrm{basis}\) say precisely that the indicated vectors are
bases of these summands.  Their union is a basis.  The proof remains
valid if \(S\) has nilpotents or is the zero ring.
\end{proof}
% END CNC2

% BEGIN CNC3
\paragraph{CNC3. Compatible formal coordinates with the specified tangent vectors.}
\label{claim:CNC3}
There exist \(d\), a commutative formal group law \(F\) over \(S\),
a formal \(O\)-action \(a\mapsto[a]\), coordinate vectors
\(u_\iota\in S^d\), and maps
\(x_m:\operatorname{Spec}J_m(S,1)\to A_S\) which satisfy every
field of the original \(\mathrm{IsChart}\) at \(y\).
\begin{proof}
Take \(d=|I|\) and enumerate the basis of CNC2.  Its dual is a basis
of the conormal module \(\mathcal I/\mathcal I^2\) of the zero
section.  The section is closed since \(A_S\to\operatorname{Spec}S\)
is separated.  Its order-\(m\) infinitesimal neighbourhood is affine:
a scheme which is a nilpotent thickening of an affine scheme is affine.
Write its ring as \(B_m\), with augmentation \(B_m\to S\).
The structural map \(S\to B_m\) splits this augmentation.

The smooth-section coordinate theorem gives, locally on the base,
\(\operatorname{gr}_{\mathcal I}\mathcal O_{A_S}
 \simeq\operatorname{Sym}_S(\mathcal I/\mathcal I^2)\).
This isomorphism is intrinsic and so glues.  Choose lifts of the given
conormal basis to the augmentation ideal of \(B_1\), then inductively
lift them through \(B_{m+1}\twoheadrightarrow B_m\).
Surjectivity on these global sections follows because the neighbourhoods
are affine and the kernels are quasi-coherent.  Subtracting the
augmentation preserves the chosen lift and places it in the augmentation
ideal.  The resulting maps
\[
 S[t_1,\ldots,t_d]/(t)^{m+1}\longrightarrow B_m              \tag{CNC.4}
\]
are isomorphisms: their maps on the degree-zero term are identities and
on every positive associated graded term are the symmetric powers of
the chosen conormal basis.  Induction on the finite filtration proves
injectivity and surjectivity of (CNC.4).  They commute with truncation.
This proves the needed global chart over the affine base, rather than
assuming that one local neighbourhood contains the whole zero section.

Transport the actual group multiplication and its action maps through
these formal coordinates.  They give series \(F(x,z)\) and \([a](x)\)
with zero constant term.  The identities of the group and action hold
at every jet order, hence coefficientwise: \(F\) is associative and
commutative with zero as identity; \([a+b](x)=F([a]x,[b]x)\),
\([ab]=[a]\circ[b]\), and \([1]=\mathrm{id}\).
The maps \(x_m\) are the inverse charts followed by the actual
inclusion into \(A_S\).  Their zero and truncation identities, their
group-law identity and their action identity are the four corresponding
fields of \(\mathrm{IsChart}\).

The differential of (CNC.4) is the chosen dual basis.  Therefore its
coordinate tangent vectors are the prescribed vectors of \(y\).
Take \(u_\iota\) to be the corresponding standard coordinate vector.
This proves the fields \(\mathrm{tan}\) and \(\mathrm{basis}\),
including uniqueness of every linear combination.  No Noetherian
assumption on \(S\) has been used.
\end{proof}
% END CNC3

% BEGIN CNC4
\paragraph{CNC4. Nilpotent lifting and normalized global primitive data.}
\label{claim:CNC4}
An invertible module over a ring with nilpotent ideal admits a generator
lifting any generator of its reduction.  The given primitive and adjoint
line-bundle isomorphisms consequently admit the normalizations used in
the original definitions \(\mathrm{IsPrim}\) and \(\mathrm{IsFrame}\).
\begin{proof}
Let \(R\) be a ring, \(J\) a nilpotent ideal and \(P\) an invertible
\(R\)-module.  Lift the prescribed generator of \(P/JP\) to \(p\in P\).
At every prime of \(R\), choose a basis of the rank-one free module
\(P_\mathfrak p\).  The coefficient of \(p\) in this basis is a
unit modulo \(J_\mathfrak p\), hence is a unit, since \(J\) is
nilpotent.  Thus \(R\to P\), \(1\mapsto p\), is an isomorphism
at every prime and therefore an isomorphism.  This proves the lifting
statement without a local-ring assumption.

Write the cube-defect line as
\[
 B_L(x,z)=L_{x+z}\otimes L_x^{-1}\otimes L_z^{-1}\otimes L_0.
                                                               \tag{CNC.5}
\]
The field \(\mathrm{SAB}.\mathrm{prim}\) is the actual isomorphism
\(L\simeq B_L(\mathrm{id},[w_0])\).  Equivalently it gives a
global generator of
\(L\otimes B_L(\mathrm{id},[w_0])^{-1}\).
At zero the second factor has its canonical trivialization.  The ratio
of this generator with the given rigidification of \(L_0\) is a
unit on the original base.  Multiply by the inverse pullback of that
unit.  The resulting generator \(w_P\) has exactly the zero-section
value required by \(\mathrm{IsPrim}\).

Similarly \(\mathrm{SAB}.\mathrm{adj}(a)\) supplies an isomorphism
between \(B_L([a]x,z)\) and \(B_L(x,[a^*]z)\).  Choose its
generator \(w_a\) in
\(B_L(x,[a^*]z)\otimes B_L([a]x,z)^{-1}\), and normalize its
value at \((0,0)\) to one.  This is possible by a base unit and
unique by CNC1.  All these choices are made for the original line
bundle on the original family.  In particular the primitive identity
is not asserted for an arbitrary symmetric line bundle lacking the
original \(\mathrm{prim}\) field.
\end{proof}
% END CNC4

% BEGIN CNC5
\paragraph{CNC5. Compatible normalized frames and their four formal ratios.}
\label{claim:CNC5}
For the chart of CNC3 there are compatible generators \(e_m\) of
the pullbacks of \(B_L\) to \(J_m(S,2)\), equal to the canonical
generator on either zero axis.  They define formal units
\(P(x,z,w),Q(x,z,w),s(x,z),a_b(x,z)\) which give all the original
\(\mathrm{IsFrame}\) equations, with \(P=\mathrm{ph1}\),
\(Q=\mathrm{ph2}\), \(s=\mathrm{sy}\), and
\(a_b=\mathrm{ad}(b)\).
\begin{proof}
Pull \(L\) back along the formal sum of the two universal points.
Starting from its rigidification at zero, repeatedly apply CNC4 to the
surjective jet truncations.  This gives compatible generators \(g_m\)
of this line at every order.  Divide by their restrictions to the two
axes and multiply by their value at their intersection.  The result is
the compatible generator
\[
 e_m(x,z)=g_m(x,z)\,g_m(x,0)^{-1}g_m(0,z)^{-1}g_m(0,0)
                                                               \tag{CNC.6}
\]
of (CNC.5).  The formula is an equality of line tensors, so it is
independent of local trivializations.  It gives one on both axes.

Let \(c(x,z,w)\) denote the generator of the triple cube defect
specified by the original cubical structure.  At each jet order take
the following ratios of generators of the same line:
\[
\begin{aligned}
 P(x,z,w)&=c(x,z,w)\,
              \frac{e(x,w)e(z,w)}{e(F(x,z),w)},\\
 Q(x,z,w)&=c(z,w,x)\,
              \frac{e(x,z)e(x,w)}{e(x,F(z,w))},\\
 s(x,z)&=\frac{e(z,x)}{e(x,z)},\\
 a_b(x,z)&=w_b(x,z)\,
                    \frac{e([b]x,z)}{e(x,[b^*]z)}.
\end{aligned}                                                   \tag{CNC.7}
\]
Here and below a fraction is a tensor product with inverse generators;
each whole expression is an actual scalar unit.  Ratios commute with
truncation, so these units define unique formal series over \(S\).
Their constant terms are one, by the axes of \(e\), the normalized
cube, and CNC4.  Conversely their truncations are exactly the units in
the source \(\mathrm{IsFrame}\) equations: multiplying (CNC.7)
by its displayed denominator recovers, in order, its first partial
law, second partial law with the permutation \((z,w,x)\), symmetry
law, and adjoint law.  This also checks the direction of the adjoint
ratio and the inverse in the source symmetry equation.
\end{proof}
% END CNC5

% BEGIN CNC6
\paragraph{CNC6. All formal biextension identities on the original frames.}
\label{claim:CNC6}
The series from CNC5, together with \(F\) and \([b]\), are a
literal object \(B:\mathrm{FB}(S,d,O)\).
\begin{proof}
For brevity in this proof only, write \(x+z=F(x,z)\) and
\(b x=[b](x)\).  The ten zero-axis conditions are
\[
 P(0,z,w)=P(x,0,w)=P(x,z,0)=1,
 \quad Q(0,z,w)=Q(x,0,w)=Q(x,z,0)=1,
\]
\[
 s(0,z)=s(x,0)=a_b(0,z)=a_b(x,0)=1.                            \tag{CNC.8}
\]
The remaining identities are the following; this list includes every
non-axis field of the original \(\mathrm{FB}\):
\[
\begin{aligned}
 P(x,z,t)P(x+z,w,t)&=P(z,w,t)P(x,z+w,t),\\
 P(z,x,w)&=P(x,z,w),\\
 Q(x,z,w)Q(x,z+w,t)&=Q(x,w,t)Q(x,z,w+t),\\
 Q(x,w,z)&=Q(x,z,w),\\
 P(x,z,w)P(x,z,t)Q(x+z,w,t)
   &=Q(x,w,t)Q(z,w,t)P(x,z,w+t),\\
 s(x,z)s(z,x)&=1,\\
 Q(x,z,w)s(x,z)s(x,w)&=s(x,z+w)P(z,w,x).
\end{aligned}                                                   \tag{CNC.9}
\]
For every \(b,c\in O\), the adjoint identities are
\[
\begin{aligned}
 P(bx,bz,w)a_b(x+z,w)
   &=a_b(x,w)a_b(z,w)P(x,z,b^*w),\\
 Q(bx,z,w)a_b(x,z+w)
   &=a_b(x,z)a_b(x,w)Q(x,b^*z,b^*w),\\
 a_1(x,z)&=1,\\
 a_{bc}(x,z)&=a_b(cx,z)a_c(x,b^*z),\\
 P(bx,cx,z)a_{b+c}(x,z)
   &=a_b(x,z)a_c(x,z)Q(x,b^*z,c^*z),\\
 a_b(x,z)a_{b^*}(z,x)s(bx,z)&=s(x,b^*z).
\end{aligned}                                                   \tag{CNC.10}
\]

We verify these as line-tensor identities before reading their
coefficients.  The cubical structure gives a symmetric normalized
trivialization of the triple cube defect.  Its cocycle equation is the
agreement of the two ways to split a sum of three points; its
interchange equation is the agreement of splitting in the first and
second variables.  Inserting (CNC.7), the \(e\)-factors cancel from
the first, third and fifth equations of (CNC.9), leaving respectively
these two cocycle equations and the interchange equation for the cube.
The second and fourth use its symmetry.  The sixth cancels immediately.
In the seventh both sides reduce to the same cube at \((z,w,x)\)
and the same ratio
\(e(z,x)e(w,x)/e(x,z+w)\).  The faces of the cube and (CNC.6)
give the first eight axis equations.

For completeness, the normalized adjoint isomorphisms used here have
all the compatibility diagrams which these cancellations require.
The two ways to combine the adjoint isomorphism with the first or
second cubical law are isomorphisms of the same line bundles on
\(A^3\).  The two composites for a product \(bc\), for a sum
\(b+c\), and for transposing the variables are isomorphisms of
the same line bundles on \(A^2\).  They agree at the zero tuple:
every cube face is one and every \(w_b\) was normalized there.
Their ratios are therefore one by CNC1.  The same argument on \(A\)
shows that either zero-axis restriction of \(w_b\) is one.  The
isomorphism for \(1\) is likewise the identity.  Substitution of
(CNC.7) into these six diagrams gives exactly (CNC.10), including its
factors \(P(bx,cx,z)\) and \(Q(x,b^*z,c^*z)\).
This proves the last two equations of (CNC.8) as well.

All line identifications just used are between tensor expressions in
the same \(L\), with the indicated group and \(O\)-action maps;
associativity, commutativity, and \([bc]=[b][c]\) identify their
source and target.  Thus no biextension compatibility of an arbitrarily
chosen adjoint isomorphism was presumed: it follows from its
normalization and CNC1.  Pull the identities to each universal jet.
Equality of every truncation is equality of formal series, so
(CNC.8)--(CNC.10) prove the asserted \(\mathrm{FB}\) structure
over the entire \(S\).
\end{proof}
% END CNC6

% BEGIN CNC7
\paragraph{CNC7. Formal logarithms over a general rational algebra.}
\label{claim:CNC7}
Let \(R\) be a commutative \(\mathbb Q\)-algebra and \(F\) a
commutative \(d\)-dimensional formal group law over \(R\).
There is a unique formal coordinate change \(\ell(x)=x+O(x^2)\)
with
\(\ell(F(x,z))=\ell(x)+\ell(z)\).  In these coordinates every
formal \(O\)-action is linear.
\begin{proof}
Set \(J(x)=\partial_zF(x,0)\), an invertible matrix with
\(J(0)=1\).  Associativity, differentiated in the third variable,
gives
\[
 \partial_zF(x,z)J(z)=J(F(x,z)).                              \tag{CNC.11}
\]
The vector fields
\(X_i=\sum_jJ_{ji}(x)\partial_{x_j}\) commute.  Indeed
commutativity and associativity give
\(F(F(x,se_i),te_j)=F(F(x,te_j),se_i)\); the coefficient of
\(st\) is exactly \([X_i,X_j]=0\).
The rows of the one-form \(\omega=J(x)^{-1}dx\) are dual to
these fields.  For each row,
\(d\omega(X_i,X_j)=X_i(\omega(X_j))-X_j(\omega(X_i))
-\omega([X_i,X_j])=0\).  Since the \(X_i\) form a basis of
derivations, \(d\omega=0\).

Every closed formal one-form over a rational algebra has its radial
primitive with zero constant term: for a row \(\sum_i b_i(x)dx_i\),
it is
\(\int_0^1\sum_i x_i b_i(tx)\,dt\).
Here integration means dividing a coefficient of \(t^m\) by
\(m+1\).  Differentiation and the equality
\(\partial_jb_i=\partial_ib_j\) verify the primitive directly.
Apply this construction to \(\omega\) to get \(d\ell=\omega\),
\(\ell(0)=0\), and linear part the identity.
Equation (CNC.11) shows that the derivative in \(z\) of
\(\ell(F(x,z))-\ell(z)\) is zero.  A formal series over a
\(\mathbb Q\)-algebra with all first derivatives zero is constant,
coefficient by coefficient.  Evaluation at \(z=0\) gives the desired
addition identity.

A coordinate series with invertible linear part has a unique inverse,
obtained one homogeneous degree at a time.  Any two logarithms differ
by an additive coordinate series with linear part the identity.
For an additive series \(q\), differentiating
\(q(x+z)=q(x)+q(z)\) at \(z=0\) shows that its derivative is
constant.  It is therefore linear.  This proves uniqueness, and the
same argument applies to
\(\ell\circ[b]\circ\ell^{-1}\), giving a matrix \(M(b)\).
The formal action laws show that \(b\mapsto M(b)\) is a unital
ring homomorphism to \(\operatorname{Mat}_d(R)\).
\end{proof}
% END CNC7

% BEGIN CNC8
\paragraph{CNC8. A normalized symmetric additive two-cocycle is a coboundary.}
\label{claim:CNC8}
Suppose \(c(x,z)\in R[[x,z]]\), with \(R\) a commutative
\(\mathbb Q\)-algebra, satisfies
\[
 c(x,0)=c(0,z)=0,\quad c(x,z)=c(z,x),\quad
 c(x,z)+c(x+z,w)=c(z,w)+c(x,z+w).
\]
Then a series \(h(x)\), with zero constant and linear parts, satisfies
\(c(x,z)=h(x)+h(z)-h(x+z)\).
The construction is linear in the coefficients of \(c\), so it also
works with formal parameter variables.
\begin{proof}
Put \(b_i(x)=\partial_{z_i}c(x,0)\).  Differentiating the cocycle
identity in \(w_i\) at zero gives
\[
 \partial_{z_i}c(x,z)=b_i(x+z)-b_i(z).                        \tag{CNC.12}
\]
By symmetry,
\(\partial_{x_j}c(x,z)=b_j(x+z)-b_j(x)\).
Equality of the two mixed derivatives therefore gives
\(\partial_jb_i=\partial_ib_j\).  Also \(b_i(0)=0\).
Define
\[
 h(x)=-\int_0^1\sum_i x_i b_i(tx)\,dt.                       \tag{CNC.13}
\]
The same elementary radial calculation as in CNC7 gives
\(\partial_i h=-b_i\).  Thus the derivative in each \(z_i\)
of \(h(x)+h(z)-h(x+z)\) is the right side of (CNC.12).
The difference from \(c\) is independent of \(z\) and vanishes
at \(z=0\); it is zero.  The origin and linear terms of \(h\)
vanish because the \(b_i\) have zero constant term.  Formula
(CNC.13) uses only derivatives, rational divisions and linear
operations.  It respects any linear conditions in other parameter
variables, including vanishing and linear dependence on those
variables.  For parameters we apply it coefficient by coefficient;
each resulting bidegree depends on finitely many coefficients.
\end{proof}
% END CNC8

% BEGIN CNC9
\paragraph{CNC9. The precise gauge change.}
\label{claim:CNC9}
For any formal \(\mathrm{FB}\) over a \(\mathbb Q\)-algebra
and any \(h(x,z)\) vanishing on both axes, define
\[
\begin{aligned}
 \delta_1h(x,z,w)&=h(x,w)+h(z,w)-h(F(x,z),w),\\
 \delta_2h(x,z,w)&=h(x,z)+h(x,w)-h(x,F(z,w)),\\
 \delta_sh(x,z)&=h(z,x)-h(x,z),\\
 \delta_bh(x,z)&=h([b]x,z)-h(x,[b^*]z).
\end{aligned}                                                   \tag{CNC.14}
\]
Replacing the four units \((P,Q,s,a_b)\) by
\[
 (P\exp(-\delta_1h),\ Q\exp(-\delta_2h),\
       s\exp(-\delta_sh),\ a_b\exp(-\delta_bh))               \tag{CNC.15}
\]
again gives a formal \(\mathrm{FB}\), with the same group law and
action.  Successive gauge changes add their \(h\)'s.
\begin{proof}
Each series in (CNC.14) has zero constant term, so its exponential is
defined.  Formal exponential is multiplicative on sums in a commutative
rational algebra, as follows from the binomial formula at each finite
degree.  Substitute (CNC.15) into (CNC.8)--(CNC.10).  The old unit
identity cancels, leaving an additive identity in \(h\).
For the first partial associativity identity this is
\[
 \delta_1h(x,z,t)+\delta_1h(F(x,z),w,t)
 =\delta_1h(z,w,t)+\delta_1h(x,F(z,w),t),
\]
where all terms cancel by associativity of \(F\).  The other
partial associativity is the same calculation in the second variable;
the two commutativities and interchange follow by permuting or
regrouping the four displayed summands.  The symmetry identities use
\(\delta_sh(z,x)=-\delta_sh(x,z)\).

The adjoint compatibility equations follow by applying
\([b]F=F([b],[b])\) in the two variables.  For the product equation
the needed cancellation is
\[
 \delta_{bc}h(x,z)=\delta_bh([c]x,z)
                         +\delta_ch(x,[b^*]z).
\]
For the sum equation use
\([b+c]x=F([b]x,[c]x)\) and the corresponding identity for the
involution.  Finally
\(\delta_bh(x,z)+\delta_{b^*}h(z,x)+\delta_sh([b]x,z)
=\delta_sh(x,[b^*]z)\), which is immediate on expanding four
differences.  This verifies every adjoint law, and the axis laws follow
from those of \(h\).  Addition of gauges follows from linearity of
all four \(\delta\)'s and multiplicativity of exponential.
\end{proof}
% END CNC9

% BEGIN CNC10
\paragraph{CNC10. Killing the two partial cocycles.}
\label{claim:CNC10}
For every formal \(\mathrm{FB}\) over a commutative
\(\mathbb Q\)-algebra there is an axis-zero gauge \(h_{12}\)
after which \(P=Q=1\).
\begin{proof}
Pass through CNC7 to additive coordinates.  Formal logarithm is inverse
to exponential between the ideal of series with zero constant term and
units with constant term one.  This follows by the usual one-variable
formal identities, whose substitutions are valid in the positive-degree
ideal.  Taking logarithms turns all unit identities into additive
identities over the same coefficient algebra.

Write \(p=\log P\).  For each formal parameter \(w\), the
first two identities in (CNC.9) say that \(p(x,z,w)\) is a
normalized symmetric two-cocycle in \(x,z\).  CNC8 constructs
\(h_1(x,w)\) with \(\delta_1h_1=p\).  It has zero first-variable
origin, and vanishes at \(w=0\) because \(p\) does.  Gauge by
\(h_1\).  The new \(P\) is one, and all identities remain true
by CNC9.

Let \(q=\log Q\) for the new system.  Interchange in (CNC.9)
now says
\(q(x+z,w,t)=q(x,w,t)+q(z,w,t)\).
Differentiating in \(z\) at zero shows that \(q\) is linear
in \(x\).  The remaining partial associativity and symmetry say
that it is a normalized symmetric two-cocycle in its last two variables.
Apply CNC8 there, coefficient by coefficient in the first variable,
to obtain \(h_2(x,z)\) with \(\delta_2h_2=q\).
It is linear in \(x\) and vanishes on both axes, so
\(\delta_1h_2=0\).  A second gauge kills \(Q\) while preserving
\(P=1\).  Thus \(h_{12}=h_1+h_2\) has the required property
in additive coordinates.  Substituting the formal logarithm back into
both variables gives the assertion in the original coordinates.
\end{proof}
% END CNC10

% BEGIN CNC11
\paragraph{CNC11. Killing every adjoint cocycle by one separability calculation.}
\label{claim:CNC11}
Suppose that \(O\) is a commutative reduced ring finite and free over
\(\mathbb Z\), with an involution.  After CNC10, one bilinear gauge
kills \(a_b\) simultaneously for every \(b\in O\), and preserves
\(P=Q=1\).
\begin{proof}
Remain in additive coordinates and write the action as the matrices
\(M(b)\) of CNC7.  The first two equations of (CNC.10), with
\(P=Q=1\), say that \(\log a_b(x,z)\) is additive in each
variable.  It is therefore bilinear, so write it as \(x^TD(b)z\).
The remaining equations give
\[
 D(b+c)=D(b)+D(c),\quad D(1)=0,\quad
 D(bc)=M(c)^TD(b)+D(c)M(b^*).                                 \tag{CNC.16}
\]
Commutativity of \(O\) permits exchanging \(b,c\) in the last
equation.  With
\(\mathcal L(b)Z=M(b)^TZ\) and
\(\mathcal R(b)Z=ZM(b^*)\), it is consequently the ordinary
bimodule derivation identity
\(D(bc)=\mathcal L(b)D(c)+D(b)\mathcal R(c)\).
The two actions commute.  The maps extend \(\mathbb Q\)-linearly
to \(B=O\otimes\mathbb Q\); this is allowed because the target
coefficient algebra is rational.  The identity extends as well.

The algebra \(B\) is finite and reduced over \(\mathbb Q\), hence
a product of finite separable fields.  Let
\(e=\sum_i u_i\otimes v_i\in B\otimes_{\mathbb Q}B\)
be its separability element.  Explicitly it is the sum of a trace-dual
basis tensor in each field factor.  Its two identities are
\[
 \sum_i u_iv_i=1,\qquad
 (b\otimes1)e=e(1\otimes b).                                 \tag{CNC.17}
\]
One can check them after a splitting-field extension: under
\(B\otimes\overline{\mathbb Q}\simeq
\overline{\mathbb Q}^{\operatorname{Hom}(B,\overline{\mathbb Q})}\),
the tensor is the diagonal idempotent, so both identities are coordinate
identities.  Faithful scalar extension brings them back to \(B\).
This also covers a zero algebra by the empty sum on its zero module.

Set \(Z=\sum_i\mathcal L(u_i)D(v_i)\).  Applying the second
identity of (CNC.17) to the linear map
\(u\otimes v\mapsto\mathcal L(u)D(v)\), and then using the
derivation identity, gives
\[
 \mathcal L(b)Z
 =\sum_i\mathcal L(u_i)D(v_ib)
 =D(b)+Z\mathcal R(b).
\]
Thus \(D(b)=M(b)^TZ-ZM(b^*)\).
The bilinear series \(h_3(x,z)=x^TZz\) has
\(\delta_bh_3=\log a_b\) for every \(b\), while its two
partial coboundaries are zero.  CNC9 therefore kills every adjoint
unit in a single gauge change and preserves the already trivial
partial cocycles.  The separability calculation takes place over
\(\mathbb Q\); it does not require the possibly nonreduced
coefficient algebra to be a field or an etale algebra.
\end{proof}
% END CNC11

% BEGIN CNC12
\paragraph{CNC12. The full original canonical-gauge theorem.}
\label{claim:CNC12}
For every commutative ring \(R_0\), every \(d\), every commutative
reduced finite free \(\mathbb Z\)-algebra \(O\) with involution,
every commutative \(\mathbb Q\)-algebra \(R\) which is also an
\(R_0\)-algebra, and every \(B:\mathrm{FB}(R_0,d,O)\), there is
\(h\in R[[x,z]]\) with
\[
 h(0,z)=h(x,0)=0,\qquad
 (P,Q,s,a_b)_R
   =(\exp\delta_1h,\exp\delta_2h,
      \exp\delta_sh,\exp\delta_bh)\quad(b\in O).              \tag{CNC.18}
\]
This is the complete original \(\mathrm{CanonGaugeSt}\), hence in
particular supplies \(B.\mathrm{IsCanonGauge}(S,h)\) for CNC6.
\begin{proof}
Map the coefficients of \(B\) to \(R\); all its equations persist.
Use CNC10 and CNC11 in additive coordinates.  At this stage
\(P=Q=a_b=1\) for every \(b\).  Put \(t=\log s\).
The last two symmetry equations of (CNC.9) give
\[
 t(x,z)=-t(z,x),\qquad t(x,z+w)=t(x,z)+t(x,w).
\]
They imply that \(t\) is bilinear.  The last equation of (CNC.10)
now gives \(t(M(b)x,z)=t(x,M(b^*)z)\).
Take \(h_4=-t/2\).  Then
\(\delta_sh_4=t\); its two partial coboundaries vanish by
bilinearity, and every \(\delta_bh_4\) vanishes by the displayed
adjoint relation.  This last gauge makes all four units equal to one.

Add the four gauges of CNC10--CNC11 and this last step, and substitute
\(\ell(x),\ell(z)\) into the resulting additive-coordinate
series.  CNC9 shows that the original units are exactly the four
exponentials in (CNC.18), with the original formal group law and
original action.  This is precisely the six fields \(\mathrm{ax0}\),
\(\mathrm{ax1}\), \(\mathrm{ph1}\), \(\mathrm{ph2}\),
\(\mathrm{sy}\), and \(\mathrm{ad}\) of
\(\mathrm{IsCanonGauge}\), including the sign of the symmetry
difference and the involution in the adjoint difference.  Since \(O_E\)
is a reduced finite free \(\mathbb Z\)-algebra, CNC6 is an allowed
instance.  Every construction is valid in all finite jet orders over
the same \(R\), so no reducedness or continuity condition has entered.
\end{proof}
% END CNC12

% BEGIN CNC13
\paragraph{CNC13. The actual line generator and the unique germ of every section.}
\label{claim:CNC13}
Let \(w_P\) and \(e_m\) be the generators of CNC4--CNC5.
For every \(k\geq0\) and every original section
\(F\in H^0(A_S,L^{\otimes k})\) there is a unique series
\(f\in S[[x_1,\ldots,x_d]]\) satisfying the original
\(\mathrm{IsGerm}\).  Its coefficient at every jet order is
evaluation of the given section in the given line generator.
\begin{proof}
Define the generator of \(L\) at the formal point by
\[
 e_{L,m}(x)=w_P(x)\,e_m(x,[w_0]x).                           \tag{CNC.19}
\]
The factors belong to \(L_x\otimes B_L(x,[w_0]x)^{-1}\)
and \(B_L(x,[w_0]x)\), respectively, so their product belongs
to the actual \(L_x\).  Their generator properties show that it is
a generator.  It commutes with truncation, by CNC3 and CNC5.
At the zero jet it equals the original rigidification, by the
normalization of \(w_P\) and the axes of \(e_m\).

Pull the section \(F\) to the actual family over each jet algebra,
then evaluate it at \(x_m\) in the basis \(e_{L,m}^{\otimes k}\).
This gives a unique element \(f_m\in J_m(S,1)\).
Compatibility of pullback, evaluation and (CNC.19) gives
\(f_{m+1}\bmod(t)^{m+1}=f_m\).
The inverse limit of these rings is the formal series ring over \(S\):
the coefficient of a monomial is its coefficient in any truncation
large enough to contain it.  Let \(f\) be this series.

The pullback section just used is the \(F'\) in the definition of
\(\mathrm{IsGerm}\), the tensor power of (CNC.19) is its
\(\tau\), and the equality \(f\bmod(t)^{m+1}=f_m\) is its
\(\mathrm{IsValueAt}\) clause.  Conversely any witnesses in that
definition have the same pullback section and the same generator under
its prescribed correspondence.  The scalar coefficient in a generator
is unique.  Thus every such series has all truncations \(f_m\),
which proves uniqueness.  This uses equality in \(S\), not equality
on its points, and therefore also proves the assertion for nonreduced
coefficient algebras.
\end{proof}
% END CNC13

% BEGIN CNC14
\paragraph{CNC14. A complete original chart at every framed point.}
\label{claim:CNC14}
For every commutative \(K\)-algebra \(S\), with its rational-algebra
structure, and every \(y:\mathrm{FrPt}(X_o,c,S)\), there exists
an original \(C:\mathrm{ChartAt}(X_o,c,y)\) with compatible frames
and a germ for every \(k\geq0\) and every original section.
\begin{proof}
Use the involution of (CNC.1).  CNC2--CNC3 supply
\(d,F,[b],u,x_m\) and all six chart fields.  CNC4 supplies the
global primitive generator \(w_P\) and its \(\mathrm{IsPrim}\)
normalization.  CNC5 supplies \(e_m\) and all four frame equations;
CNC6 packages their series as the original formal \(\mathrm{FB}\).
CNC12 supplies its gauge \(h\).  Formula (CNC.19) at zero is the
field \(\mathrm{rig0}\), by CNC13.  These are all the data and
all the proof fields in \(\mathrm{ChartAt}\).  Frame truncation,
although not separately demanded in that structure, is available from
CNC5 and is exactly \(\mathrm{FrTruncC}\) used by the source
germ producer.  Finally CNC13 supplies, for every natural \(k\) and
every section, the required germ.

The result thus matches the entire conclusion of
\texttt{ChartRing.exists\_chartAt\_with\_germs} on the original
moduli family.  The large source \(\mathrm{Laws}\) package is not
being taken as a new premise: of its geometric content this construction
needed only the base-functions theorem of CNC1, whose hypotheses were
checked there.  The smooth formal chart and all gauge equations have
been constructed above.
\end{proof}
% END CNC14

% BEGIN CNC15
\paragraph{CNC15. Construction of the literal original \(\mathrm{IsNab}\).}
\label{claim:CNC15}
For every original \(E,H,D_w,l,X_o,c\) there exist the star-ring
structure \(a^*=H.\mathrm{bar}(a)\) on \(O_E\) and
\(\nabla:\mathrm{NabV}(X_o,c)\) such that
\(\mathrm{IsNab}(X_o,c,\nabla)\).
This has the full coefficient-algebra and section quantifiers of
\texttt{GlobalTheory.CanonConn.exists\_isNab\_uncond}.
\begin{proof}
For every coefficient algebra and framed point, choose once a chart
\(C\) supplied by CNC14.  For its formal group law define the
invariant differential operator in direction \(v\in S^d\) by
\[
 D_v f=\sum_{j=1}^d\left(\sum_{i=1}^d
       v_i\,\partial_{z_i}F_j(x,0)\right)\partial_{x_j}f.
                                                               \tag{CNC.20}
\]
For a multi-index \(\alpha:I\to\mathbb N\) of finite support,
take the product of \(D_{u_\iota}^{\alpha(\iota)}\) in the
original finite-index order.  This is exactly
\(\mathrm{DpolyOf}(F,u,X^\alpha)\) in the source.  The operators
also commute, by the commuting invariant fields in CNC7, so the value
would be independent of the order.

For \(k\geq0\), let \(f\) be the unique germ of \(F\) from
CNC13 and set
\[
 \nabla_y(k,F)(\alpha)=
 \left[\prod_{\iota\in I}D_{u_\iota}^{\alpha(\iota)}
    \left(f(x)\exp\{k\,h(x,[w_0]x)\}\right)\right]_{x=0}.
                                                               \tag{CNC.21}
\]
This is an element of \(S\); every requested coefficient involves
only finitely many coefficients of \(f,F,h\).
For a negative integer \(k\), set the corresponding value to zero.
This defines the function on every integer required by
\(\mathrm{NabV}\); \(\mathrm{IsNab}\) constrains only its
nonnegative values.

A \(K\)-algebra has its canonical rational-algebra structure, since
\(K\) is a number field.  Any other rational-algebra structure appearing
in the \(\mathrm{IsNab}\) quantifier has the same structure map:
a unital ring map from \(\mathbb Q\) is determined by the images
of integers and their inverses.  Thus the choice of chart and the value
in (CNC.21) can be made for the \(\mathrm{NabV}\) data alone,
and work for every such additionally specified rational structure.
The star equation is true by construction.  For every \(S,y\), the
chosen chart has all the required germs by CNC14.  If \(f'\) is
\emph{any} series satisfying its \(\mathrm{IsGerm}\), CNC13
gives \(f'=f\).  Formula (CNC.21) is then exactly
\[
 C.B.\mathrm{nabVal}(S,C.h,H.w_0,k,f',C.uvec,X^\alpha),
\]
because the source \(\mathrm{hL}\) is \(h(x,[w_0]x)\) and
its \(\mathrm{nabVal}\) is the constant coefficient after
(CNC.20).  This proves the equality for every germ, not merely a selected
representative.  It is the final field of \(\mathrm{IsNab}\),
so all its quantifiers and fields have been verified.
\end{proof}
% END CNC15

% BEGIN CNC16
\paragraph{CNC16. Exact removal of the connection premise from the lattice route.}
\label{claim:CNC16}
On every supplied original geometric family to which ICL11--ICL21
apply, CNC15 supplies their original \(\mathrm{IsNab}\) producer
at every row and for both original Boolean level choices.  The
conditional lattice-spanning implications of that chapter consequently
require no separate connection-existence premise.
\begin{proof}
CNC15 quantifies over every \(D_w,l,X_o,c\), without a rank,
weight, prime, or embedding exclusion.  Apply it to each row's actual
\(Q.X_o(j)\), with that row's original coefficient field and level.
Choosing these connections gives exactly the data assumed at the start
of ICL11; its identity cover and all subsequent displayed maps use those
same objects.  If two Boolean choices use different original levels,
apply the same universally quantified theorem to each.  No change in
the \(\mathrm{FWTie}\) maps of ICL21 is needed.

This supplies one producer in that conditional route.  CNC15 asserts
neither construction of the common moduli family nor integral
valuation bounds for (CNC.21).  In particular the principal
same-Gram charts, actual finite theta system, Gaussian identity and
absolute integrability, global short-rule dictionary, and ramified
short-value estimates retained by ICL16--ICL21 are still required.
The higher-cover and all-tuple class-sum obligations are not consequences
of \(\mathrm{IsNab}\) alone.  Keeping those hypotheses gives exactly
the stated implication and makes no additional whole-obligation closure.
\end{proof}
% END CNC16
